\documentclass{zhvt-classic}

% 使用本地源流明体（GenRyuMin2 TW）按文件路径加载，不依赖系统字体名
% 版框（内框／界栏／书口侧栏／鱼尾）改用与版心印章同一印泥棕色：
%   印章图片 yshs.png 实色 = RGB(135,68,52) = HTML 874434
\usepackage{xcolor}
\usepackage{tikz}
% sealbrown 双注册：xcolor 供 TikZ/pgf，L3 供 cls 版框
\definecolor{sealbrown}{HTML}{874434}
\ExplSyntaxOn
\color_set:nnn { sealbrown } { HTML } { 874434 }
\ExplSyntaxOff

\zhvtset{
  grid color       = sealbrown,
  main font        = GenRyuMin2TW-SB,
  main font path   = C:/Users/Administrator/texmf/fonts/opentype/public/GenRyuMin2TW/,
jiazhu font      = GenRyuMin2TW-M,
  jiazhu font path = C:/Users/Administrator/texmf/fonts/opentype/public/GenRyuMin2TW/,
  % 页码框容纳 4 字：正常页「九十七」+前置空 1 格=4 行恰好；
  % 100+ 页「一百零一」共 5 行超出，末字「一/二/三」被框裁掉（2026-10-04 实测），故放宽到 5
  shukou page number chars = 5,
  % 书口书名框高从默认 6 字减到 4 字，把省出的 2 字高拨给 chapter 框，
  % 使 6 字书名+卷名放进 9 字高的 chapter 框后不再压线（2026-10-05）
  shukou title chars = 4,
}

\title{史记}[汲郑列传]
\author{汉·司马迁}
\maker{闲叟漫记·重製}


% ===== 索隐述赞样式（2026-11-08 四次定稿 · 愚大人亲手改定参数）=====
% ⚠ 竖排前提：自定义字体家族必须带 RawFeature={vertical:+vert}；
%   simhei.ttf 加 +vhal 会整体加载失败静默回退，只写 +vert；simfang 同理。
\newCJKfontfamily{\shuzanfont}[Path=C:/Windows/Fonts/,RawFeature={vertical:+vert}]{simhei.ttf}
\newCJKfontfamily{\shuzanbody}[Path=C:/Windows/Fonts/,RawFeature={vertical:+vert}]{simfang.ttf}
% 题头：152x30pt 圆角 6pt 棕底白字；baseline (0,2) 使整框落入界行格子不压版框线；
% 白字节点 x 不补偿、y +8.2pt（像素级墨迹实测 dx/dy≈0，L/R 留白约 20pt）。
\newcommand{\shuzanhead}[1]{%
  \begin{tikzpicture}[baseline={(0,2)}]
    \fill[sealbrown, rounded corners=6pt]
      (-76pt,-15pt) rectangle (76pt,15pt);
    \node[text=white, inner sep=0pt, anchor=center] at (0pt,8.2pt)
      {\shuzanfont\fontsize{23}{23}\selectfont #1};
  \end{tikzpicture}%
}
% ★ 修正书口卷名位置：chapter 框起始位置改回 cls 默认值，避免全局影响。
%   只在每个卷首页通过 \juan 命令里的空命令微调该页位置（安全，不会跨页污染）。
\ExplSyntaxOn
\AtBeginDocument{
  \dim_set:Nn \l__zhvt_shukou_chapter_title_start_dim
    { \l__zhvt_shukou_title_height_dim + \l__zhvt_font_size_dim * 2 }
  \dim_set:Nn \l__zhvt_shukou_page_number_start_dim
    { \l__zhvt_shukou_chapter_title_start_dim + \l__zhvt_shukou_chapter_title_height_dim }
}
\ExplSyntaxOff

\begin{document}

% ★★ 版心（书口）排布覆写：把书名放到鱼尾之下（2026-10-04）
%
%   cls 默认版心自上而下：书名（title 静态框，在鱼尾之上）→ 鱼尾 → 篇名（chapter 动态框）→ 页码
%   本模板改为：          鱼尾（位置不动）→ 书名 → 篇名 → 页码分割线 → 页码
%
%   手法：把 cls 放在 title 静态框里的书名，改放进 chapter 动态框（鱼尾之下）；
%         title 静态框留空。书名取 \l__zhvt_title_str，即 \title{...} 的第一参数，
%         所以换书时只改 \title，不必在此重复写书名。
%
%   ⚠ 三条必须遵守的写法（2026-10-04 实测）：
%     ① 汉字必须用 \hbox{...} 包住。flowfram 会把框内容装进一个宽 14.4mm（= 一个 line_sep
%        = 书口栏宽）的纵向 minipage 并按该宽度折行，而整页靠 PDF /Rotate 90 显示、
%        框内文字本身不变向 —— 不包 \hbox 时「围炉夜话」4 字会被逐字折成 4 行，
%        在版面上就成了横排 4 字并横向溢出版框。
%     ② 竖排下**一个 \hbox 就是一整列**。书名与篇名要同处一列，必须放进同一个 \hbox；
%        分成两个 \hbox 时，后一个会被当成新的一列、甩到隔壁栏去
%        （实测「序」跑到书口栏外侧 14.4mm 的正文列上）。中间的 \quad 在框内是横向空白，
%        显示出来正好是书名与篇名之间的纵向间隔。
%     ③ ★ 内容盒要写成 **与框等宽的定宽盒** `\hbox to \l__zhvt_shukou_chapter_title_height_dim{…\hfil}`。
%        实测：chapter 框是把内容盒沿纵向**居中**摆放的，内容宽度一变，书名就跟着漂：
%        自然宽度的 \hbox（书名 38.4mm）落在 x=157.2mm，加篇名后（57.6mm）落在 x=152.4mm，
%        目录页（无篇名）与正文页的书名就差半格。做成等宽盒后，无论有没有篇名，
%        内容固定从框顶起排，实测两页均落在 x=142.8mm，不再漂。
\ExplSyntaxOn
% 把书口 chapter 框（书名+篇名）起始位置下移 1 字高，避免 4 字卷名顶到鱼尾
% ⚠ cls 的 \__zhvt_calculate_dimensions: 在 \AtBeginDocument 里执行，会重算该 dim；
%   必须把覆写挂到 begin document 之后，否则会被二次覆盖（2026-10-05 实测）。
% (旧方案已移到导言区，直接修改 chapter 框起始位置)
% 卷标题命令 #1=书口篇名(鱼尾) #2=正文卷名（含卷序，卷名内空格用 ~）
%
% ★★ 目录页两项修订（2026-10-05，愚大人指出）★★
%   ① 页码起点统一对齐：
%      原状 —— 目录项写成 \par\hspace{4em}卷名\hspace{4em}页码，页码起点 = 卷名实宽 + 4em。
%      卷名实宽随字数变：「卷一~梁惠王上」6 字 +~、「卷七~离娄上」5 字 +~，
%      差 1 个字位 → 卷七~卷十 4 条的页码整体比其余 10 条高 27.21pt（一个字），看着上下不齐。
%      改法 —— 卷名补全角空格（宽 = 1 字 = 27.215pt）补足到 **7 字符位**再写进目录，
%      14 条页码起点遂全部落在同一条竖线上（起点对齐；末字随页码字数自然参差）。
%      ⚠ 只对目录里的副本 \l_tmpc_tl 补，正文卷名仍用 \l_tmpa_tl 原样 —— 正文像素不受影响。
%   ② 卷名整体上移半字位：前导 \hspace 由 4em 改 3.5em。
%      实测原状目录块在版框内偏下（天头留白 139.6pt / 地脚留白 110.7pt），
%      上移 0.5em = 13.6pt 后上下基本均分，视觉居中。
\NewDocumentCommand { \juan } { mm } {
	\clearpage
	\tl_set:Nn \l__zhvt_chapter_title_tl { #1 }
	\tl_set:Nn \l_tmpa_tl { #2 }
	\tl_set_eq:NN \l_tmpc_tl \l_tmpa_tl
	\int_step_inline:nn { 7 - \str_count:n { #2 } }
	  { \tl_put_right:Nn \l_tmpc_tl { 　 } }
	\tl_set:Nx \l_tmpb_tl { \zhnum { page } }
	\tl_put_left:Nn \l_tmpb_tl { \hspace { 4em } }
	\tl_put_left:Nn \l_tmpb_tl { \l_tmpc_tl }
	\tl_put_left:Nn \l_tmpb_tl { \par \hspace { 3.5em } }
	\addtocontents { toc } { \l_tmpb_tl }
	\skip_horizontal:n { \l__zhvt_font_size_dim * 4 }
	\l_tmpa_tl
	\par
	\thispagestyle { empty }
}
\cs_set_protected:Npn \__zhvt_set_shukou_content: {
	% ★ 鱼尾上方留白（title 静态框）嵌入闲叟钤印：
	%   整页 /Rotate 90，故此处把图用 \rotatebox{90} 转正为竖章；
	%   尺寸卡在版心条幅内（书内约 7mm 宽 × 33.8mm 高，留余量不压鱼尾）。
	\setstaticcontents * { title } {
		\vbox_to_ht:nn { \l__zhvt_baseline_skip_dim } {
			\vskip 14.2pt
			\hbox_to_wd:nn { \l__zhvt_shukou_title_height_dim } {
				\hfil
				\rotatebox [ origin = c ] { 90 } { \includegraphics [ width = 7mm , height = 33.8mm ] { yshs.png } }
				\hfil
			}
		}
	}
% ★ 版心（书口）统一改法（2026-11-16 愚大人定）：18pt＋框内居中＋书口文字同版框色
%   ① 原 27.2pt 时 7 字长篇名（如「廉颇蔺相如列传」）末字溢出 chapter 框、压住页码；
%      改 \fontsize{18pt}{18pt}，前后 \hfil 令字串在框内居中（不随篇名字数漂移）。
%   ② \vbox{\vskip 8.2pt} 定量下移，抵消缩字号后栏宽方向 8.2pt 上偏（像素实测 y 中心 543.3 = 原版）。
%      ⚠ \raisebox / \vspace* 在 flowfram 框内无效；换字号须同步换补偿：16pt→10.1、19pt→7.4。
%   ③ chapter 与 page~number 均置 \color{sealbrown}（HTML 874434），与版框／界行同色；正文仍黑。
%   ⚠ 页码字位、钤印、鱼尾、界行、正文一律未动；五部合订本同此口径。
	\setdynamiccontents * { chapter }
	{ \color{sealbrown} \vbox { \vskip 8.2pt \hbox to \l__zhvt_shukou_chapter_title_height_dim
		{ \hfil \fontsize{18pt}{18pt}\selectfont \l__zhvt_title_str \quad \l__zhvt_chapter_title_tl \hfil } } }
	\setdynamiccontents * { page~number } { \color{sealbrown} \quad \zhnum { page } }
}
\__zhvt_set_shukou_content:
\ExplSyntaxOff

\maketitle

% 单篇，无多卷目录，不渲染目录页

\juan{汲郑列传}{卷一百二十~汲郑列传第六十}
\hspace*{2em}汲黯字长孺。濮阳人也。其先有宠於古之卫君。至黯七世。世为卿大夫。黯以父任。孝景时为太子洗马。以庄见惮。孝景帝崩。太子即位。黯为谒者。东越相攻。上使黯往视之。不至。至吴而还。报曰。﹁越人相攻。固其俗然。不足以辱天子之使。﹂河内失火。延烧千馀家。上使黯往视之。还报曰。﹁家人失火。屋比延烧。不足忧也。臣过河南。河南贫人伤水旱万馀家。或父子相食。臣谨以便宜。持节发河南仓粟以振贫民。臣请归节。伏矫制之罪。﹂上贤而释之。迁为荥阳令。黯耻为令。病归田里。上闻。乃召拜为中大夫。以数切谏。不得久留内。迁为东海太守。黯学黄老之言。治官理民。好清静。择丞史而任之。其治。责大指而已。不苛小。黯多病。卧闺閤内不出。岁馀。东海大治。称之。上闻。召以为主爵都尉。列於九卿。治务在无为而已。弘大体。不拘文法。\\
\hspace*{2em}黯为人性倨。少礼。面折。不能容人之过。合己者善待之。不合己者不能忍见。士亦以此不附焉。然好学。游侠。任气节。内行脩絜。好直谏。数犯主之颜色。常慕傅柏。袁盎之为人也。善灌夫。郑当时及宗正刘弃。亦以数直谏。不得久居位。\\
\hspace*{2em}当是时。太后弟武安侯蚡为丞相。中二千石来拜谒。蚡不为礼。然黯见蚡未尝拜。常揖之。天子方招文学儒者。上曰吾欲云云。黯对曰。﹁陛下内多欲而外施仁义。柰何欲效唐虞之治乎。﹂上默然。怒。变色而罢朝。公卿皆为黯惧。上退。谓左右曰。﹁甚矣。汲黯之戆也。﹂群臣或数黯。黯曰。﹁天子置公卿辅弼之臣。宁令从谀承意。陷主於不义乎。且已在其位。纵爱身。柰辱朝廷何。﹂\\
\hspace*{2em}黯多病。病且满三月。上常赐告者数。终不愈。最後病。庄助为请告。上曰。﹁汲黯何如人哉。﹂助曰。﹁使黯任职居官。无以逾人。然至其辅少主。守城深坚。招之不来。麾之不去。虽自谓贲育亦不能夺之矣。﹂上曰。﹁然。古有社稷之臣。至如黯。近之矣。﹂\\
\hspace*{2em}大将军青侍中。上踞厕而视之。丞相弘燕见。上或时不冠。至如黯见。上不冠不见也。上尝坐武帐中。黯前奏事。上不冠。望见黯。避帐中。使人可其奏。其见敬礼如此。\\
\hspace*{2em}张汤方以更定律令为廷尉。黯数质责汤於上前。曰。﹁公为正卿。上不能襃先帝之功业。下不能抑天下之邪心。安国富民。使囹圄空虚。二者无一焉。非苦就行。放析就功。何乃取高皇帝约束纷更之为。公以此无种矣。﹂黯时与汤论议。汤辩常在文深小苛。黯伉厉守高不能屈。忿发骂曰。﹁天下谓刀笔吏不可以为公卿。果然。必汤也。令天下重足而立。侧目而视矣。﹂\\
\hspace*{2em}是时。汉方征匈奴。招怀四夷。黯务少事。乘上间。常言与胡和亲。无起兵。上方向儒术。尊公孙弘。及事益多。吏民巧弄。上分别文法。汤等数奏决谳以幸。而黯常毁儒。面触弘等徒怀诈饰智以阿人主取容。而刀笔吏专深文巧诋。陷人於罪。使不得反其真。以胜为功。上愈益贵弘。汤。弘。汤深心疾黯。唯天子亦不说也。欲诛之以事。弘为丞相。乃言上曰。﹁右内史界部中多贵人宗室。难治。非素重臣不能任。请徙黯为右内史。﹂为右内史数岁。官事不废。\\
\hspace*{2em}大将军青既益尊。姊为皇后。然黯与亢礼。人或说黯曰。﹁自天子欲群臣下大将军。大将军尊重益贵。君不可以不拜。﹂黯曰。﹁夫以大将军有揖客。反不重邪。﹂大将军闻。愈贤黯。数请问国家朝廷所疑。遇黯过於平生。\\
\hspace*{2em}淮南王谋反。惮黯。曰。﹁好直谏。守节死义。难惑以非。至如说丞相弘。如发蒙振落耳。﹂\\
\hspace*{2em}天子既数征匈奴有功。黯之言益不用。\\
\hspace*{2em}始黯列为九卿。而公孙弘。张汤为小吏。及弘。汤稍益贵。与黯同位。黯又非毁弘。汤等。已而弘至丞相。封为侯。汤至御史大夫。故黯时丞相史皆与黯同列。或尊用过之。黯褊心。不能无少望。见上。前言曰。﹁陛下用群臣如积薪耳。後来者居上。﹂上默然。有间黯罢。上曰。﹁人果不可以无学。观黯之言也日益甚。﹂\\
\hspace*{2em}居无何。匈奴浑邪王率众来降。汉发车二万乘。县官无钱。从民贳马。民或匿马。马不具。上怒。欲斩长安令。黯曰。﹁长安令无罪。独斩黯。民乃肯出马。且匈奴畔其主而降汉。汉徐以县次传之。何至令天下骚动。罢弊中国而以事夷狄之人乎。﹂上默然。及浑邪至。贾人与市者。坐当死者五百馀人。黯请间。见高门。曰。﹁夫匈奴攻当路塞。绝和亲。中国兴兵诛之。死伤者不可胜计。而费以巨万百数。臣愚以为陛下得胡人。皆以为奴婢以赐从军死事者家。所卤获。因予之。以谢天下之苦。塞百姓之心。今纵不能。浑邪率数万之众来降。虚府库赏赐。发良民侍养。譬若奉骄子。愚民安知市买长安中物而文吏绳以为阑出财物于边关乎。陛下纵不能得匈奴之资以谢天下。又以微文杀无知者五百馀人。是所谓﹃庇其叶而伤其枝﹄者也。臣窃为陛下不取也。﹂上默然。不许。曰。﹁吾久不闻汲黯之言。今又复妄发矣。﹂後数月。黯坐小法。会赦免官。於是黯隐於田园。\\
\hspace*{2em}居数年。会更五铢钱。民多盗铸钱。楚地尤甚。上以为淮阳。楚地之郊。乃召拜黯为淮阳太守。黯伏谢不受印。诏数彊予。然後奉诏。诏召见黯。黯为上泣曰。﹁臣自以为填沟壑。不复见陛下。不意陛下复收用之。臣常有狗马病。力不能任郡事。臣原为中郎。出入禁闼。补过拾遗。臣之原也。﹂上曰。﹁君薄淮阳邪。吾今召君矣。顾淮阳吏民不相得。吾徒得君之重。卧而治之。﹂黯既辞行。过大行李息。曰。﹁黯弃居郡。不得与朝廷议也。然御史大夫张汤智足以拒谏。言足以饰非。务巧佞之语。辩数之辞。非肯正为天下言。专阿主意。主意所不欲。因而毁之。主意所欲。因而誉之。好兴事。舞文法。内怀诈以御主心。外挟贼吏以为威重。公列九卿。不早言之。公与之俱受其僇矣。﹂息畏汤。终不敢言。黯居郡如故治。淮阳政清。後张汤果败。上闻黯与息言。抵息罪。令黯以诸侯相秩居淮阳。七岁而卒。\\
\hspace*{2em}卒後。上以黯故。官其弟汲仁至九卿。子汲偃至诸侯相。黯姑姊子司马安亦少与黯为太子洗马。安文深巧善宦。官四至九卿。以河南太守卒。昆弟以安故。同时至二千石者十人。濮阳段宏始事盖侯信。信任宏。宏亦再至九卿。然卫人仕者皆严惮汲黯。出其下。\\
\hspace*{2em}郑当时者。字庄。陈人也。其先郑君尝为项籍将。籍死。已而属汉。高祖令诸故项籍臣名籍。郑君独不奉诏。诏尽拜名籍者为大夫。而逐郑君。郑君死孝文时。\\
\hspace*{2em}郑庄以任侠自喜。脱张羽於戹。声闻梁楚之间。孝景时。为太子舍人。每五日洗沐。常置驿马安诸郊。存诸故人。请谢宾客。夜以继日。至其明旦。常恐不遍。庄好黄老之言。其慕长者如恐不见。年少官薄。然其游知交皆其大父行。天下有名之士也。武帝立。庄稍迁为鲁中尉。济南太守。江都相。至九卿为右内史。以武安侯魏其时议。贬秩为詹事。迁为大农令。\\
\hspace*{2em}庄为太史。诫门下。﹁客至。无贵贱无留门者。﹂执宾主之礼。以其贵下人。庄廉。又不治其产业。仰奉赐以给诸公。然其餽遗人。不过算器食。每朝。候上之间。说未尝不言天下之长者。其推毂士及官属丞史。诚有味其言之也。常引以为贤於己。未尝名吏。与官属言。若恐伤之。闻人之善言。进之上。唯恐後。山东士诸公以此翕然称郑庄。\\
\hspace*{2em}郑庄使视决河。自请治行五日。上曰。﹁吾闻﹃郑庄行。千里不赍粮﹄。请治行者何也。﹂然郑庄在朝。常趋和承意。不敢甚引当否。及晚节。汉征匈奴。招四夷。天下费多。财用益匮。庄任人宾客为大农僦人。多逋负。司马安为淮阳太守。发其事。庄以此陷罪。赎为庶人。顷之。守长史。上以为老。以庄为汝南太守。数岁。以官卒。\\
\hspace*{2em}郑庄。汲黯始列为九卿。廉。内行脩絜。此两人中废。家贫。宾客益落。及居郡。卒後家无馀赀财。庄兄弟子孙以庄故。至二千石六七人焉。\\
\hspace*{2em}太史公曰。夫以汲。郑之贤。有势则宾客十倍。无势则否。况众人乎。下邽翟公有言。始翟公为廷尉。宾客阗门。及废。门外可设雀罗。翟公复为廷尉。宾客欲往。翟公乃大署其门曰。﹁一死一生。乃知交情。一贫一富。乃知交态。一贵一贱。交情乃见。﹂汲。郑亦云。悲夫。\\

% ===== 索隐述赞（司马贞《史记索隐》）· 统一样式 =====
% 题头「索隐述赞曰」五字并入反白框；正文每四句一列，定宽盒 \hfil 居列中；
% \fontsize{24}{41} 列距对齐全书界行（121+41k）；\clearpage 使述赞独占一页。
\clearpage
\shuzanhead{索隐述赞曰}\linebreak
{\shuzanbody\fontsize{24}{41}\selectfont
\hbox to \hsize{\hfil 河南矫制。自古称贤。淮南卧理。天子伏焉。\hfil}\linebreak
\hbox to \hsize{\hfil 积薪兴叹。伉直愈坚。郑庄推士。天下翕然。\hfil}\linebreak
\hbox to \hsize{\hfil 交道势利。翟公怆旃。\hfil}%
\par}
\end{document}
